\chapter{The complete reading grid}
\label{chap:A-grille}

This appendix reproduces the reading grid of the survey, criterion by criterion, with the question
put to the reader and its decision rule. It combines the reader's sheet (\file{study/GRILLE.md}) and
the rules of the protocol (\file{study/PROTOCOL.md}): deviations 1 to 6 for the harmonisation of
grid v1, deviations 11, 14, 17 and 19 for grid v2. When in doubt, the protocol prevails.

\section{What the reader does}
\label{sec:a-lecteur}

The reader reads the whole paper, appendices included, before filling in one JSON file per paper
(\cref{sec:formats}). For each criterion, the file gives the value, an exact quotation from the
paper, copied word for word from the PDF, then the page. If nothing justifies a \texttt{yes}, the
quotation is \enquote{absent après lecture des sections expériences, annexes et liens} (absent after
reading the experiments, appendices and links sections) and the page is \texttt{null}. A quotation
spread over two pages receives a list of pages. In a two-column PDF, a quotation that runs across
lines is split into one-line pieces joined by \enquote{[...]}. The reader never sees another
reader's codes, the results or the synthesis document. They run \file{study/check\_codes.py} on
their own files and fix them until it reports zero problems. The possible values are
\texttt{yes}, \texttt{no} and \texttt{na} (not applicable), except for the integer criteria and for
the family X5.

\section{Pineau's grid (Q1 to Q8)}
\label{sec:a-pineau}

Joelle Pineau presented these eight criteria at NeurIPS 2018 (slides 29 and 30), with the share of
the 50 RL papers from 2018 she had read that met each one. Her sampling method is not described. The
labels below are hers, word for word.

\begin{table}[htbp]
\centering
\small
\caption{Pineau's eight criteria and her 2018 figures.}
\label{tab:a-pineau-2018}
\begin{tabular}{@{}llr@{}}
\toprule
Code & Pineau's label & 2018 \\
\midrule
Q1 & Paper has experiments & 100\% \\
Q2 & Paper uses neural networks & 90\% \\
Q3 & All hyperparams for proposed algorithm are provided & 90\% \\
Q4 & All hyperparams for baselines are provided & 60\% \\
Q5 & Code is linked & 55\% \\
Q6 & Method for choosing hyperparams is specified & 20\% \\
Q7 & Evaluations on some variation of a hold-out test set & 10\% \\
Q8 & Significance testing applied & 5\% \\
\bottomrule
\end{tabular}
\end{table}

\begin{small}
\begin{longtable}{@{}p{0.07\linewidth}p{0.27\linewidth}p{0.6\linewidth}@{}}
\caption{Grid v1, Pineau's criteria: question and decision rule.}\label{tab:a-grille-q}\\
\toprule
Code & Question & Decision rule \\
\midrule
\endfirsthead
\toprule
Code & Question & Decision rule \\
\midrule
\endhead
\bottomrule
\endfoot
Q1 & Does the paper have experiments? & \texttt{no}: every other criterion is \texttt{na} and
  leaves the denominators (deviation 3). Pineau's denominators are the experimental papers. \\
Q2 & Does it use neural networks? & Read directly. \\
Q3 & Are all the hyperparameters of the method given? & \texttt{yes} if a table or a paragraph gives
  them for all the main methods. \enquote{Default values of [library]} without listing them counts
  \texttt{yes} only if the version is given. \\
Q4 & Are all the hyperparameters of the baselines given? & Same rule as Q3. \texttt{na} if no
  baseline is compared. \\
Q5 & Is the code linked? & A link to code in the paper. \enquote{On request}, \enquote{will be
  released}, \enquote{will be made public at <URL>}: \texttt{no}. A project page counts \texttt{yes}
  only if it leads to code, checked by visiting it on 9 November 2026. An anonymous mirror that holds
  the code counts \texttt{yes} (deviations 6 and 19). \\
Q6 & Is the way the hyperparameters are chosen described? & The selection method is stated: grid,
  random search, budget, criterion. \enquote{Tuned} and nothing more: \texttt{no}. \\
Q7 & Does evaluation run on held-out conditions? & Tasks, environments, environment random seeds or
  datasets seen neither in training nor in tuning. For an \texttt{llm} paper, benchmarks distinct
  from the RL training data count \texttt{yes}. A \texttt{theory} paper with toy experiments:
  \texttt{no} (deviation 4). New training random seeds, used to train, are not held-out conditions
  (deviation 17). \\
Q8 & Is a statistical test applied to a comparison? & A test (t, Welch, Mann-Whitney, bootstrap with
  interval, rliable interquartile mean with stratified intervals) applied to a comparison. Error bars
  alone: \texttt{no} for Q8, \texttt{yes} for X2. \\
\end{longtable}
\end{small}

\section{Our additions in grid v1 (X1 to X7)}
\label{sec:a-x}

These criteria come from the ICLR 2018 talk and from the 2025 debate on the spurious rewards of
language models (\cref{sec:llm-pourquoi}). They were written with the protocol, on 8 November 2026,
before any reading.

\begin{small}
\begin{longtable}{@{}p{0.07\linewidth}p{0.27\linewidth}p{0.6\linewidth}@{}}
\caption{Grid v1, our additions: question and decision rule.}\label{tab:a-grille-x}\\
\toprule
Code & Question & Decision rule \\
\midrule
\endfirsthead
\toprule
Code & Question & Decision rule \\
\midrule
\endhead
\bottomrule
\endfoot
X1 & How many independent training runs behind the main results? & An integer, or \texttt{null} if
  not stated: the smallest number announced for a main table or curve. For an \texttt{llm} paper,
  count the independent RL training runs alone: pass@k and \enquote{16 samples} are evaluation
  samples, which do not count. \\
X2 & Is an uncertainty reported on the main results? & Standard deviation, standard error, confidence
  interval, band around the curve. \\
X3 & Are the best runs selected? & \enquote{Top-k}, \enquote{best of N}, best checkpoint chosen on
  the test set. \texttt{yes} also when the paper's own tables show that the main result is the best
  of several identical training runs, even if the text does not say so (deviation 5). \\
X4 & Does the paper say that the baselines are tuned with the same care as the method? &
  \texttt{yes} requires a tuning procedure described for the baselines, or an equal tuning budget
  announced. \enquote{Same configuration as the method}, with no tuning described for anyone:
  \texttt{no} (deviation 1). \\
X5 & Which family does the paper belong to? & \texttt{llm} if the optimised policy is a pretrained
  generative model conditioned on language (a language model of any size, a vision-language model, a
  text-to-image model). \texttt{classic} if RL trains something else, even in the service of a large
  model, or a transformer trained from scratch on RL tasks. \texttt{theory} for a paper with no
  experiments or with toy experiments only. A paper with theoretical results that compares methods on
  standard benchmarks, even in an appendix, is \texttt{classic} (deviation 18). \\
X6 & For \texttt{llm}: how many distinct base-model families? & An integer, or \texttt{null}. Qwen,
  Llama, Mistral, Gemma, DeepSeek, OLMo\dots{} Two sizes of Qwen count as one family. \\
X7 & Is the compute used quantified? & \texttt{yes} requires at least one quantity: hours, GPU-hours,
  FLOP, or a number of GPUs with a duration. The name of the hardware alone: \texttt{no}
  (deviation 2). \\
\end{longtable}
\end{small}

\section{Grid v2 (Y1 to Y12)}
\label{sec:a-y}

Grid v2 was written on 9 November 2026, before any v2 reading (deviation 11), together with its
decision rules (deviation 14). The rules were made more precise after a pilot on three papers
(deviation 17), then harmonised after the audit of the raw readings, before any percentage
(deviation 19). Each criterion points to the 2018 source that motivates it.

\paragraph{Scope rules} If Q1 is \texttt{no}, Y1 to Y12 are \texttt{na}. Y2 and Y11 are \texttt{na}
if and only if X2 is \texttt{no}. Y4 is \texttt{na} if Y3 is 1. Y5 is \texttt{na} if and only if Q4
is \texttt{na}. For a paper already read with grid v1, the reader receives the v1 values of Q1, Q4
and X2, which set these rules, without seeing the other v1 codes.

\begin{small}
\begin{longtable}{@{}p{0.07\linewidth}p{0.27\linewidth}p{0.6\linewidth}@{}}
\caption{Grid v2: question, source and decision rule.}\label{tab:a-grille-y}\\
\toprule
Code & Question & Decision rule \\
\midrule
\endfirsthead
\toprule
Code & Question & Decision rule \\
\midrule
\endhead
\bottomrule
\endfoot
Y1 & Is the performance measure defined? (ICLR 2018, slide 27) & \texttt{yes} if the paper states
  the quantity and the point of measurement: final return, mean over the last N episodes, best point,
  area under the curve, pass@k with k and the number of samples, success rate and over what.
  \enquote{Return} or \enquote{accuracy} and nothing more: \texttt{no}. A normalised score with no
  statement of the point at which it is taken: \texttt{no}. The number of evaluation episodes is not
  required (deviation 17). A curve whose evaluation points are stated counts as a point of
  measurement. \enquote{Trained for N steps} followed by one figure counts as the final
  point. For a quantity other than the score of a trained policy (error of an estimator, compute
  time), \texttt{yes} if the quantity is defined (deviation 19). \\
Y2 & Are the error bars defined? (checklist v1.2) & \texttt{yes} if the paper states what the band
  or the interval represents: standard deviation, standard error, 95\% confidence interval, minimum
  and maximum. Bars present but undefined: \texttt{no}. \\
Y3 & How many tasks does the main conclusion cover? (Henderson et al., table 3) & An integer: the
  number of distinct training setups on which the method is trained and compared, an environment or
  a training dataset for \texttt{llm}. Two difficulty levels of one game, or the quality variants of
  a dataset from one environment (D4RL medium, replay, expert), count as one task. Several evaluation
  benchmarks of one trained model do not multiply the tasks (deviation 17). A single meta-training
  counts as one task. A setup present only in an appendix does not count (deviation 19). Derived
  value for the analysis: \texttt{yes} if the integer is at least 2. \\
Y4 & Is a task where the method does not win shown? (Henderson et al., table 3) & \texttt{yes} if a
  task-by-task result is given and at least one task where the method does not win is shown or
  stated. \texttt{no} if the paper gives only means, or shows only tasks where the method wins. \\
Y5 & Is the implementation of the baselines identified? (Henderson et al., figure 6) & \texttt{yes}
  if, for each main baseline, the paper says where the code comes from: a named repository, library
  or software framework, with or without version, \enquote{reimplemented by us}, figures copied from
  papers with their source, a released model designated by its exact name, the authors' training
  framework if the baselines are said to run in it (deviations 17 and 19). A single main baseline
  with no origin gives \texttt{no}. \\
Y6 & Are exclusions stated? (checklist v1.2) & \texttt{yes} if the paper says which runs, random
  seeds or data were removed, or says explicitly that none was. An omitted combination with its
  reason counts (deviation 17). A filtering of the training data stated with what it removes counts
  too (deviation 19). Silence is \texttt{no}. \\
Y7 & Is the version of the environments or benchmarks given? (checklist v1.2) & \texttt{yes} if the
  version of each main environment or benchmark is given: Gym or Gymnasium suffix, suite version,
  exact name and date of the test split for \texttt{llm}. \enquote{MuJoCo} alone: \texttt{no}.
  \enquote{MATH-500} counts \texttt{yes}, \enquote{MATH} alone \texttt{no}. A dataset published on a
  named repository counts as versioned, a split built by the authors only if it is published
  (deviation 17). A name that covers several sets without saying which (\enquote{AMC},
  \enquote{Minerva}) does not count. An environment fully specified in the paper counts
  (deviation 19). \\
Y8 & Is a generalisation test outside the training conditions carried out? (NeurIPS 2018, roadmap)
  & \texttt{yes} if at least one evaluation runs on conditions absent from training and tuning: other
  environment random seeds announced as such, new tasks, modified dynamics, a distinct test
  distribution. Stricter than Q7: for \texttt{llm}, a different benchmark counts only if it is
  presented as an out-of-distribution test. A test split from the same distribution falls under Q7.
  Y8 \texttt{yes} presupposes Q7 \texttt{yes} (deviation 19). \\
Y9 & Is the sensitivity to at least one hyperparameter shown? (ICLR 2018, slides 24 to 26) &
  \texttt{yes} if at least one hyperparameter of the method, presented as such, is swept and the
  result is shown through a performance measure. A component ablation does not count. A parameter
  set at inference counts if it is presented as a hyperparameter. A condition of the environment does
  not count (deviations 17 and 19). \\
Y10 & Are the software dependencies listed with their versions? (checklist v1.2) & \texttt{yes} if
  the versions of the main libraries are in the paper, or if an environment file is cited in the
  linked repository. The repository is not visited, unless the paper points explicitly to the file. \\
Y11 & Is the type of uncertainty named? (comparable to Gorsane et al., 2022) & \texttt{yes} or
  \texttt{no}, plus a \texttt{kind} field: \texttt{std}, \texttt{se}, \texttt{ci}, \texttt{iqr},
  \texttt{minmax} or \texttt{other}. Broader than Y2: it is enough that the type is named somewhere.
  Percentile bands other than the interquartile range: \texttt{other} (deviation 17). \\
Y12 & Is a statistical tool from after 2018 used? (adoption of Agarwal et al., 2021) & \texttt{yes}
  if the paper uses at least one of these tools: interquartile mean, stratified bootstrap interval,
  rliable, performance profiles, probability of improvement. A plain bootstrap on a mean counts if
  it produces a reported interval. \\
\end{longtable}
\end{small}

\section{Harmonisation}
\label{sec:a-harmonisation}

A harmonisation rule is written into the protocol, dated, before the first percentage it affects. It
applies to both years, from the quotations, in \file{study/overrides.json} (grid v1) or
\file{study/overrides\_v2.json} (grid v2), each modified value with its reason. Agreement between
readers is computed on the raw readings, before harmonisation. For X3, X4 and X7, no rule was added
after the enlarged reading, so that both years are still read with the same grid (deviation 19). The
function \texttt{pineau\_rl.grid.validate} checks the schema and the scope rules of every file
(\cref{sec:formats}).
